\section{Costing method}\label{sec:method}

The first-pass costings use three public inputs and a small set of
labelled scenario parameters. Everything is generated by a single
pipeline (\path{analysis_taa/costing.py}) that downloads and pins its raw
files and emits every manuscript value as a machine-generated macro; no
reported number is typed by hand.

\paragraph{Eligible flows.} Annual redundancy flows come from the ONS
Labour Force Survey redundancy series: the four quarters to June 2026 give
\TaaFlowManuf\ manufacturing redundancies and \TaaFlowAll\ across all
industries. These are \emph{flows of events}, the right base for an
insurance costing, and they embed no trade attribution: the narrow rule
prices insurance for all manufacturing redundancy, trade-caused or not,
which is what an administrable scheme without certification would cover.

\paragraph{Earnings.} Pre-displacement earnings are anchored to ASHE
full-time weekly earnings for manufacturing
(\pounds\TaaMedianWeeklyManuf\ median, \pounds\TaaMeanWeeklyManuf\ mean)
and the whole economy (\pounds\TaaMedianWeeklyAll\ median). Re-employment earnings apply a penalty of
15 or 28 per cent (central 21.5)---the hours-controlled and
all-employee penalties calibrated on FRS data in the companion study,
imported here as scenario values rather than re-estimated.

\paragraph{Scenario parameters.} Re-employment within a year:
50/70/85 per cent, bracketing the UK
displacement literature's range \citep{hijzen2010, upwardwright2019}.
Take-up of a new benefit: bracketed low/central/high per arm, anchored to
observed contributory-benefit and UC take-up where analogues exist and
labelled pure assumptions where they do not (the retraining arm).
Contribution-condition satisfaction for the UI arm: 90 per cent central
(80/100 low/high). Each parameter enters the
pipeline once, and every cost cell is reported with its scenario triple.

\paragraph{What this method deliberately is not.} The costings are gross,
static, annual, steady-state envelopes. They do not net off the taxes and
taper clawback that reduce what payments cost the Exchequer on net
(microsimulation stage); they hold behaviour fixed, so the
\citet{hyman2024} employment response---which made US wage insurance
plausibly self-financing---enters only as a labelled sensitivity that
scales non-employment durations; they do not model the macroeconomic
state-dependence of redundancy flows, which in a downturn would be
multiples of the recent average; and the redundancy-disregard arm is
bounded, not costed. These are the same envelope conventions used by the
first-pass costings in the UK unemployment-insurance debate
\citep{ifs2025ui}, and they are the appropriate honesty for a paper whose
purpose is to put a designed instrument set on the table.
